% Arsitektur jaringan Hyperledger Fabric sistem Rupiah Digital (TikZ).
% Purwarupa studi empat organisasi, kanal bi-coin.
% Di-input dari BAB III §3.2.2 setelah tab:jaringan.

\begin{figure}[H]
\centering
\resizebox{\textwidth}{!}{%
\begin{tikzpicture}[
  peer/.style={draw=black, fill=white, rectangle, rounded corners=2pt,
    minimum width=3.2cm, minimum height=0.9cm, align=center, font=\scriptsize},
  observer/.style={draw=black, dashed, fill=none, rectangle, rounded corners=2pt,
    minimum width=3.2cm, minimum height=0.9cm, align=center, font=\scriptsize},
  cdb/.style={draw=black, fill=gray!12, rectangle,
    minimum width=2.2cm, minimum height=0.6cm, align=center, font=\scriptsize},
  ordnode/.style={draw=black, fill=black!8, rectangle, rounded corners=3pt,
    minimum width=3.4cm, minimum height=0.9cm, align=center, font=\footnotesize\bfseries},
  backend/.style={draw=black, fill=black!5, rectangle, rounded corners=3pt,
    minimum width=4.0cm, minimum height=0.9cm, align=center, font=\footnotesize},
  chanbox/.style={draw=black, dashed, rounded corners=5pt,
    fill=none, inner sep=8pt},
  ar/.style={-{Stealth}, thick},
  gar/.style={-{Stealth}, thick, dashed},
  pathlabel/.style={font=\scriptsize, midway, align=center},
]

%--- Backend + Orderer ---
\node[backend] (be)  at (-2.5, 3.5) {Backend API\\Go + Fabric Gateway SDK};
\node[ordnode] (ord) at ( 5.25, 3.5) {orderer.centralbank.bi-coin:7050\\(Raft, satu node)};
\draw[ar] (be.east) -- node[pathlabel,above]{gRPC} (ord.west);

%--- Peer nodes (empat organisasi) ---
\node[peer]     (pcb)   at ( 0.0, 1.2) {peer0.centralbank.bi-coin\\CentralBankMSP\\:17051};
\node[peer]     (pdist) at ( 3.5, 1.2) {peer0.distributor.bi-coin\\DistributorMSP\\:18051};
\node[observer] (pval)  at ( 7.0, 1.2) {peer0.validator.bi-coin\\ValidatorMSP (observer)\\:19051};
\node[peer]     (pusr)  at (10.5, 1.2) {peer0.enduser.bi-coin\\EndUserMSP\\:20051};

%--- Orderer delivers block to all peers ---
\draw[ar] (ord.south) to[out=-100,in=90] (pcb.north);
\draw[ar] (ord.south) to[out=-100,in=90] (pdist.north);
\draw[ar] (ord.south) -- (pval.north);
\draw[ar] (ord.south) to[out=-80,in=90]  (pusr.north);

%--- Backend gateway ke peer0.distributor (gateway peer contoh pada network/README.md) ---
\draw[gar] (be.south) -- node[pathlabel,above,sloped]{\textit{endorse}} (pdist.north);

%--- CouchDB state databases (satu per organisasi) ---
\node[cdb] (ccb)   at ( 0.0, -0.2) {CouchDB :17584};
\node[cdb] (cdist) at ( 3.5, -0.2) {CouchDB :18584};
\node[cdb] (cval)  at ( 7.0, -0.2) {CouchDB :19584};
\node[cdb] (cusr)  at (10.5, -0.2) {CouchDB :20584};

\draw[ar] (pcb)   -- (ccb);
\draw[ar] (pdist) -- (cdist);
\draw[ar] (pval)  -- (cval);
\draw[ar] (pusr)  -- (cusr);

%--- bi-coin bounding box ---
\node[chanbox, fit=(pcb)(pdist)(pval)(pusr)(ccb)(cdist)(cval)(cusr),
  label={[font=\footnotesize\bfseries]below:%
    \textit{bi-coin} (kebijakan endorsement per kontrak, TLS)}] {};

\end{tikzpicture}%
}
\caption{Arsitektur jaringan purwarupa studi Rupiah Digital pada kanal \texttt{bi-coin} dengan empat organisasi}
\label{fig:arsitektur-jaringan}
\end{figure}
